\scanpage{64}
\mspage{20}
Mais ces opérations de bouchage sont relativement superficielles, en cela qu'elles ne permettent pas de relier les $\mathcal T_{g,\nu}$, pour $g$ fixé, à des $\mathcal T_{g',\nu}$ avec $g'<g$. Pour ceci, il faut l'opération \inserted{\textcircled{$5^\circ$}} de gommage des cercles de découpage (= ex-cercles de recollement !).

Néanmoins, il y a, \deleted{une} \inserted{en} liaison avec ça, une autre structure importante à expliciter, c'est celle des \deleted{\unsure{une|des}} groupoïdes fondamentaux \inserted{$\Pi_1(X;R)$} des $(X,[\unsure{\widetilde{\Sigma}|\Sigma}],R)$, p\textsuperscript{ts}-base dans $R$, \inserted{pour $R\ne\phi$,} en tant que groupoïdes dépendant fonctoriellement bien sûr de $R$, et \deleted{la donnée} comme nouvelle opération) pour tte \deleted{\unsure{reconstitution}} \inserted{représentation} du groupoïde par ens finis, \inserted{[\unsure{définissant un}} d'un revêtement \inserted{$X'$} \unsure{de|des} $X$), d'un nouvel objet du \deleted{$\unsure{\mathcal T_{\mathrm{top}}}$} groupoïde $S\mathcal T_{\mathrm{top}}^\pm$. Déjà dans le cas de $S_2\mathcal T_{0,3}$ la question se pose et est non triviale --- en fait, c'est sans doute là une exigence trop ambitieuse, car

\margin{$\unsure{\widetilde{\Sigma}|\Sigma}$ ici est formellement \deleted{\unsure{vide}} vide, autant \unsure{bien|donc} \unsure{pour les|pour se} mettre dans les \unsure{cas} i.e. travailler dans $S\mathcal T_{\mathrm{top}}^\pm$ au lieu de $S\mathcal TD_{\mathrm{top}}^\pm$}
\mspage{63}
